% ============================================================
%  Chapter 1 — Introduction
% ============================================================

\chapter{Introduction}
\label{chap:intro}

\section{The modelling object}
\label{sec:intro-object}

A mortgage is a small dynamical system.  In any month it occupies one
of a small number of \emph{states} --- it is current on its payments,
or 30, 60, or 90-plus days past due, or it has terminated through
prepayment or a credit event --- and between consecutive months it
moves between these states with probabilities that depend on the
loan's characteristics and the economic environment.  The cashflows
of any mortgage portfolio or mortgage-backed security are completely
determined by these movements: a prepayment returns principal early
(a risk to the investor when rates have fallen), while the path from
delinquency through foreclosure destroys value.  Modelling the
\emph{conditional transition distribution} is therefore the natural
primitive for both credit and prepayment risk.

Formally, let $S_{i,t}\in\mathcal{S}$ denote the state of loan $i$ in
month $t$, and let $x_{i,t}\in\mathbb{R}^{d}$ collect everything
observable about the loan and its environment at $t$.  The object of
interest is
\begin{equation}
\label{eq:object}
p_\theta\!\left(j \mid x_{i,t}\right) \;\approx\;
\mathbb{P}\!\left(S_{i,t+1}=j \,\middle|\, x_{i,t}\right),
\qquad j\in\mathcal{S},
\end{equation}
a multinomial classification of next month's state.  Every model in
this dissertation --- from a frequency-count transition matrix to an
ensemble of deep neural networks --- is an estimator of
\eqref{eq:object}; they differ only in how flexibly they allow the
probabilities to depend on the covariates.

The central empirical claim of \citet{sirignano2021}, which this
dissertation tests on a different and cleaner dataset, is that the
dependence of \eqref{eq:object} on $x$ is strongly \emph{nonlinear}
and features economically meaningful \emph{interactions} between
covariates.  Linear models misspecify this dependence, and the
misspecification compounds when loan-level probabilities are
aggregated to pool-level cashflows --- the level at which
mortgage-backed securities are priced.

Agency mortgage-backed securities are among the largest fixed-income markets in the world, and their value rests entirely on how the underlying loans move between paying, delinquency, default, and prepayment. For the investor, mispricing the speed of prepayment misstates both the timing of returned principal and the value of the embedded call; for the lender and the regulator, mispricing the path into default misstates capital and provisions. The stakes of the conditional transition distribution in \eqref{eq:object} are therefore not incidental to the modelling --- they are why a more faithful estimator is worth the effort.

This study sits at the confluence of two literatures. The first is the option-theoretic tradition in mortgage modelling, which treats prepayment and default as the exercise of embedded options driven by borrower incentives \citep{schwartz1989, deng2000}. The second is the turn to flexible machine learning in credit and asset pricing: that such models sharpen consumer credit-risk forecasts \citep{khandani2010}, that their flexibility carries distributional consequences in mortgage markets \citep{fuster2022}, and that nonlinear methods materially improve the measurement of expected returns where linear models fail \citep{gukellyxiu2020}. \citet{sirignano2021} bring these together for mortgage risk specifically, and it is their central claim --- that transition risk depends on covariates nonlinearly and interactively --- that this dissertation tests on a cleaner, fully reproducible dataset.

Two questions organise the work. Does the deep-learning advantage documented on subprime-heavy, private-label data survive in the prime conforming credit box, where credit events are rarer and relationships plausibly tamer? And across which regimes --- the COVID-era forbearance episode, the 2020--21 refinancing wave, the 2022--23 rate shock --- does that advantage widen, compress, or invert?

\section{Design of the study}
\label{sec:intro-design}

The analysis is organised as a horse race between estimators of
\eqref{eq:object} of increasing flexibility, in three phases.  First,
an exploratory phase establishes the dataset's scope and produces
\emph{model-free} evidence of nonlinearity and covariate interactions
(Chapter~\ref{chap:data}).  Second, loan-level models are estimated
and compared strictly out of sample under a rolling, expanding-window
design: an empirical transition matrix, multinomial logistic
regression, deep neural networks of varying depth, and an ensemble
(Chapters~\ref{chap:methodology} and~\ref{chap:loanlevel}).  Third, the
fitted monthly models are composed into twelve-month horizon
predictions and compared at the pool level --- predicted versus
realised prepayment and delinquency counts, and their translation
into valuation units (Chapter~\ref{chap:poollevel}).

These questions are answered on the Fannie Mae Single-Family Loan Performance dataset --- 3.31 billion loan-month observations across 57.6 million loans, 2000--2025 --- under a rolling, strictly out-of-sample design spanning the eleven test years 2015--2025.

\section{Contributions}
\label{sec:intro-contributions}

The dissertation deliberately adopts a replication \emph{design} ---
the cleanest way to attribute findings to data and regime rather than
to modelling idiosyncrasies --- but four elements are original.

\begin{enumerate}
\item \textbf{A prime, conforming testbed.}  The evidence in
  \citet{sirignano2021} comes from private-label data with a large
  subprime component, 1995--2014.  This study asks whether the
  deep-learning advantage survives in the clean GSE credit box, where
  events are rarer and relationships plausibly tamer.  An early answer
  is already in hand: the architecture search at full scale selects
  \emph{three} hidden layers over the five preferred in the paper ---
  the nonlinearity premium is real but structurally simpler in prime
  credit.
\item \textbf{Eleven test years the literature has not examined,}
  2015--2025, including the COVID-era forbearance episode, the
  2020--21 refinancing wave, and the 2022--23 rate shock --- regimes
  that post-date the paper's sample entirely.
\item \textbf{A stricter evaluation protocol whose output is itself a
  finding.}  Replacing the paper's single temporal split with an
  eleven-window rolling backtest produces a \emph{time series of the
  value of nonlinearity}: in which market regimes does the deep
  model's edge widen, compress, or invert?  The single-split design
  cannot ask this question.
\item \textbf{An economic translation for the pool level.}  Predicted
  and realised pool prepayments are converted into each pool's
  prepayment-speed path, run through a pass-through cashflow engine,
  and compared as errors in cashflow timing and price --- restating
  the model comparison in the units an investor in mortgage pools
  actually uses, rather than in log-likelihood alone.
\end{enumerate}

\section{Outline}
\label{sec:intro-outline}

Chapter~\ref{chap:methodology} develops the methodology: the
transition model and its estimation at scale, the rolling
out-of-sample design, the model hierarchy, the evaluation criteria,
and the pool-level framework including the economic translation.
Chapter~\ref{chap:data} describes the dataset, the construction of
the estimation panel, the macroeconomic covariates, and the
exploratory evidence that motivates nonlinear modelling.
Chapter~\ref{chap:loanlevel} reports the loan-level results, and
Chapter~\ref{chap:poollevel} the pool-level valuation.
Chapter~\ref{chap:conclusions} concludes.
